\section{Computational Study}\label{sec:experiments}

This section evaluates the computational performance of the proposed exact solution
framework and the effect of outsourcing on the resulting production plans. The
first experiment benchmarks the proposed continuous-time NG-DSSR pricing algorithm
against the time-indexed pricing framework of \citet{bulhoes2020exact}. We then
compare the two formulations of the outsourcing decision and examine the effects of
outsourcing availability and quantity discounts. All algorithms were implemented in
Java and use IBM ILOG CPLEX 22.1.1 as the LP solver. Each run uses a single thread
and is subject to a time limit of 10,800 seconds.
% Add the final server CPU, memory, and operating system before reporting results.

\subsection{Instance Generation}\label{sec:instance-generation}

We derive the test instances from the single-machine total weighted
earliness--tardiness instances used by \citet{tanaka2009exact} and distributed in
Tanaka's online repository at
\url{https://sites.google.com/site/shunjitanaka/sips/benchmark-results-sips}.
The instances generated for this study are available at
\textit{[data repository URL]}. The source files specify a processing time, a due
date, and earliness and tardiness weights for each job. By retaining or extracting
jobs from the source instances, we construct five fixed job sets for each
$n\in\{20,40,50,60,80,100\}$, giving 30 base job sets in total.

For each base job set, we construct three parallel-machine variants, with
$m\in\{2,3,4\}$ for $n\leq 60$ and $m\in\{3,4,5\}$ otherwise. Because the source
due dates are defined on a single-machine time horizon and some job sets are
extracted from larger instances, we adjust them for both the retained workload and
parallel processing. Let $\mathcal J^0$ denote the source job set and
$\mathcal J\subseteq\mathcal J^0$ the retained job set. We define the retained
workload share and the due-window center as
\[
r=\frac{\sum_{i\in\mathcal J}p_i}{\sum_{i\in\mathcal J^0}p_i},
\qquad
c_j=\left\lceil\frac{\operatorname{round}(d_j^0r)}{m^{\mathrm{ref}}}\right\rceil,
\]
where $d_j^0$ is the source due date. The reference machine counts
$m^{\mathrm{ref}}=3$ for $n\leq60$ and $m^{\mathrm{ref}}=4$ otherwise are the
middle values of the corresponding machine sets. Using a common reference value
keeps the due windows unchanged across the three actual machine levels, so that
varying $m$ changes internal capacity rather than the delivery requirements. The
center $c_j$ determines the location of the due window, while its half-width $W_j$
determines the allowable interval around that location. We set $W_j=0$ in the
zero-width class and independently draw $W_j$ from $\{80,\ldots,120\}$ and
$\{280,\ldots,320\}$ in the narrow and wide classes, respectively. Hence,
\[
[d_j^e,d_j^l]=[\max\{0,c_j-W_j\},c_j+W_j].
\]

For every job set, we generate random and family sequence-dependent setup matrices.
The random setups contain no explicit job grouping and are converted into a directed
metric. The family setups are constructed from the same base metric by partitioning
the jobs into three to six approximately balanced families and adding a changeover
component to between-family transitions, reflecting the smaller within-family and
larger between-family setups commonly used in family scheduling
\citep{schutten1996parallel,schaller2000flowline}. Both setup structures have an
average setup time of approximately $0.5\bar p$, where $\bar p$ is the average
processing time, no setup time exceeds $\bar p$, and the directed triangle
inequality is satisfied. In the family instances, the average between-family setup
is approximately four to five times the average within-family setup. This common
setup intensity lies between the small and large setup classes used by
\citet{bulhoes2020exact}. Setup costs are defined as $\kappa_{ij}=20s_{ij}$.

We complete each original-scale production environment by specifying the outsourcing
tariff. Following scheduling studies in which outsourcing costs scale with processing
requirements \citep{lee2011twostage}, we use $p_j$ as the workload basis and adjust
the quotation for the timing importance of each job. The standard quotation is
\[
q_j=p_j\max\{\alpha_j,\beta_j\}.
\]
We obtain three outsourcing price levels by setting the undiscounted quotation to
$b_j(\lambda)=\lambda q_j$, where $\lambda\in\{0.5,1,2\}$. Unless stated otherwise,
the tariff applies a continuous incremental quantity discount, consistent with the
structures studied in production outsourcing and private-fleet routing
\citep{lu2021discount,dabia2019rich}. To use a common absolute tariff across all
instance sizes, the two breakpoints $Q_1$ and $Q_2$ are fixed at 25\% and 50\% of
the median total standard quotation among the 50-job sets. The same breakpoints are
then used for every instance. The marginal rates in the three intervals are $1$,
$0.85$, and $0.70$, respectively.

The 30 base job sets, three machine levels, two setup structures, and three
due-window classes yield 540 original-scale production environments. The three
outsourcing price levels expand this set to 1,620 complete instances. This common
instance set supports both the formulation comparison and the subsequent sensitivity
analysis, while the internal-pricing experiment uses the same production data after
removing the outsourcing option.

\subsection{Comparison of Internal Pricing Methods}\label{sec:pricing-comparison}

The first experiment compares the proposed NG-DSSR pricing algorithm with the
time-indexed baseline of \citet{bulhoes2020exact}, considered both without and with
rank-1 cuts. Outsourcing is disabled so that the comparison is driven only by the
internal scheduling problem. Because the size of a time-indexed graph depends
directly on the magnitude of the time parameters,
we supplement the original scale with medium- and high-scale variants. For each job,
the processing time and due-window center are independently multiplied by an integer drawn from
$\{5,\ldots,15\}$ in the medium-scale class and $\{15,\ldots,25\}$ in the high-scale
class. Setup times and due-window widths are adjusted to the corresponding scale:
each setup matrix follows the realized aggregate processing-time multiplier, while
the narrow and wide half-width ranges are multiplied by 10 and 20, respectively.
The earliness and tardiness weights remain unchanged. This design yields 1,620
scheduling instances and 4,860 runs. The results compare the three configurations
across job sizes, machine levels, due-window classes, setup structures, and time
scales, thereby identifying the instance conditions under which each representation
is computationally most suitable.

\subsection{Comparison of Outsourcing Formulations}\label{sec:outsourcing-formulation-comparison}

The second experiment examines how the representation of outsourcing affects exact
solution performance. We compare formulations \textup{(SP1)} and \textup{(SP2)} on
the 1,620 complete instances defined in Section~\ref{sec:instance-generation}. Both
formulations use the same NG-DSSR method for internal pricing and the same initial
solution for each instance. Consequently, the comparison isolates a single modelling
choice: outsourced job sets are either generated as columns or represented explicitly
in the restricted master problem. The two formulations produce 3,240 runs in total.
Their performance is compared across job sizes, machine levels, due-window classes,
setup structures, and outsourcing price levels, with the results reported in the
tables below.

\subsection{Sensitivity Analysis}\label{sec:sensitivity-analysis}

After selecting the better-performing outsourcing formulation, we turn from
computational performance to the structure of the optimal production plan. The two
analyses below separate the effect of making outsourcing available from the effect
of offering quantity discounts. They reuse the paired instances and solutions from
the preceding experiments whenever possible.

\subsubsection{Impact of Outsourcing}\label{sec:impact-outsourcing}

This analysis assesses both the value of making outsourcing available and the
response of the optimal plan to its price. For each original-scale production
environment, we compare the solution without outsourcing with the solutions under
the low, medium, and high price levels. The analysis focuses on changes in the use
of outsourcing, the composition of total cost, and the structure of the remaining
internal schedule under different levels of machine availability, due-window width,
and setup structure. It therefore captures both the direct outsourcing payment and
the scheduling benefit that may arise when difficult jobs are removed from internal
production.

\subsubsection{Impact of Quantity Discounts}\label{sec:impact-discounts}

Following the comparative design of \citet{dabia2019rich}, we finally examine whether
a quantity discount merely lowers the payment for jobs that would already be
outsourced or changes the outsourcing decision itself. We use the five 50-job sets
at the original time scale and the medium outsourcing price. For each combination
of three machine levels, three due-window classes, and two setup structures, the
default marginal rates $(1,0.85,0.70)$ are compared with a tariff that applies the
undiscounted rate throughout, giving 90 paired production environments. The analysis
then examines whether the discount changes the extent and composition of outsourcing
and, consequently, the organization and cost of internal production.
